Real-world multi-agent systems---from robotic swarms to autonomous vehicle fleets---often involve agents operating at fundamentally different timescales. A natural question is whether such temporal heterogeneity spontaneously gives rise to hierarchical coordination strategies, where slow agents learn to implicitly model and compensate for fast agents' delayed feedback. We investigate this question by training Independent PPO agents with heterogeneous action frequencies (1$\times$, 2$\times$, and 4$\times$) across three cooperative gridworld environments that vary in their synchronization demands. We measure emergent hierarchy through cross-agent mutual information between value functions at multiple time lags. Across 45 experimental runs (5 seeds $\times$ 3 environments $\times$ 3 conditions), we find \emph{no evidence} of emergent hierarchical credit assignment: mutual information between agent value functions shows no directional asymmetry from slow to fast agents. Instead, temporal heterogeneity produces task-dependent performance effects---a 45\% advantage in complementary-role tasks ($p < 0.0001$, $d = 6.56$) but a 17\% disadvantage in synchronization-dependent tasks ($p = 0.001$, $d = -4.83$). Temporal context features significantly improve coordination only when synchronization is critical ($p = 0.042$, $d = 1.53$). These results suggest that hierarchical temporal credit assignment must be explicitly designed into multi-agent architectures rather than expected to emerge from temporal diversity alone.
